\chapter{CƠ SỞ LÝ THUYẾT VÀ CÔNG NGHỆ SỬ DỤNG}

\section{Công nghệ Frontend}
\subsection{React 18+}
React là thư viện JavaScript mã nguồn mở được phát triển bởi Facebook để xây dựng giao diện người dùng. Dự án sử dụng React 18+ với các tính năng hiện đại:

\begin{itemize}
    \item \textbf{Functional Components:} Sử dụng function-based components thay vì class components (cấu trúc hiện đại)
    \item \textbf{React Hooks:} useState, useEffect, useContext, useReducer để quản lý state và side effects
    \item \textbf{Concurrent Rendering:} Khả năng render không đồng bộ để cải thiện hiệu suất
    \item \textbf{Suspense \& Code Splitting:} Tối ưu hóa tải ứng dụng
\end{itemize}

\subsection{TypeScript}
Dự án sử dụng TypeScript với \textbf{strict mode} được bật để:

\begin{itemize}
    \item Cung cấp type safety - phát hiện lỗi tại thời điểm phát triển
    \item Tìm kiếm nhanh hơn (IDE intellisense)
    \item Dễ dàng tái cấu trúc code lớn
    \item Giảm bugs trong production
\end{itemize}

\subsection{Tailwind CSS}
Framework CSS utility-first cho phép:

\begin{itemize}
    \item Xây dựng giao diện responsive nhanh chóng
    \item Tùy chỉnh design mà không viết custom CSS
    \item Giảm kích thước file CSS (tree-shaking unused styles)
    \item Đảm bảo consistency trong design tokens
\end{itemize}

\subsection{Quản lý state}
Ứng dụng sử dụng \textbf{React Context API} kết hợp \textbf{useReducer} để:

\begin{itemize}
    \item Quản lý state toàn cục (global state)
    \item Chia sẻ data giữa các components mà không prop drilling
    \item Cung cấp authentication context, user profile, notifications
\end{itemize}

\section{Công nghệ Backend}

\subsection{ASP.NET Core 8.0}
Framework web hiện đại của Microsoft cho xây dựng RESTful API:

\begin{itemize}
    \item \textbf{Dependency Injection:} Built-in DI container
    \item \textbf{Middleware Pipeline:} Xử lý requests theo trình tự
    \item \textbf{CORS Support:} Cho phép cross-origin requests từ frontend
    \item \textbf{Built-in Logging:} Logging và diagnostics
\end{itemize}

\subsection{Entity Framework Core 8.0}
ORM (Object-Relational Mapper) mạnh mẽ cung cấp:

\begin{itemize}
    \item \textbf{Code-First Approach:} Định nghĩa database schema qua C\# code
    \item \textbf{Migrations:} Quản lý version schema database
    \item \textbf{LINQ Queries:} Viết query bằng C\# thay vì SQL
    \item \textbf{Lazy Loading \& Eager Loading:} Tối ưu hóa số lượng queries
    \item \textbf{Relationships:} Quản lý One-to-Many, Many-to-Many relations
\end{itemize}

\subsection{ASP.NET Core Identity}
Hệ thống quản lý người dùng tích hợp:

\begin{itemize}
    \item Quản lý registration, login, password reset
    \item Role-based access control (RBAC)
    \item Claims-based authorization
    \item Phân quyền (Users, Admins, Moderators)
\end{itemize}

\subsection{JWT (JSON Web Tokens)}
Token-based authentication cho API:

\begin{itemize}
    \item Stateless authentication (không cần session server)
    \item Token chứa claims về user (id, roles, permissions)
    \item Signature bảo mật token không bị giả mạo
    \item Expire time giới hạn thời gian sử dụng token
\end{itemize}

\subsection{SignalR}
Thư viện real-time communication của ASP.NET Core:

\begin{itemize}
    \item \textbf{WebSocket Support:} Kết nối two-way real-time
    \item \textbf{Automatic Fallback:} Fallback sang Server-Sent Events nếu WebSocket không available
    \item \textbf{Hub Pattern:} Quản lý server-client messaging
    \item \textbf{Broadcasting:} Gửi notifications cho nhiều clients cùng lúc
\end{itemize}

\section{Cơ sở dữ liệu}

\subsection{SQL Server}
Hệ quản trị cơ sở dữ liệu quan hệ của Microsoft:

\begin{itemize}
    \item \textbf{Scalability:} Hỗ trợ cơ sở dữ liệu lớn
    \item \textbf{Reliability:} ACID transactions, backup \& recovery
    \item \textbf{Indexes:} Tối ưu hóa query performance
    \item \textbf{Constraints:} Đảm bảo data integrity
\end{itemize}

\subsection{Entity Framework Core Integration}
Giao tiếp qua Entity Framework Core bằng cơ chế \textbf{Code-First}:

\begin{itemize}
    \item Định nghĩa entities như C\# classes
    \item EF Core tự generate migrations
    \item `Update-Database` command cập nhật schema
    \item DbContext quản lý DbSets cho mỗi entity
\end{itemize}

\section{Công nghệ khác}

\subsection{Swagger/OpenAPI}
API documentation tự động tạo từ code:

\begin{itemize}
    \item Interactive API explorer (Swagger UI)
    \item Tự động generate từ annotations/comments
    \item Giúp frontend developers hiểu API endpoints
\end{itemize}

\subsection{Azure Services}
Triển khai trên cloud:

\begin{itemize}
    \item \textbf{Azure App Service:} Host cho Web API
    \item \textbf{Azure SQL Database:} Cơ sở dữ liệu cloud
    \item \textbf{Azure Blob Storage:} Lưu trữ hình ảnh/files
    \item \textbf{Azure Key Vault:} Bảo mật connection strings \& secrets
    \item \textbf{Application Insights:} Monitoring \& diagnostics
\end{itemize}

\subsection{GitHub \& Version Control}
Quản lý source code:

\begin{itemize}
    \item \textbf{Git:} Version control system
    \item \textbf{GitHub:} Repository hosting \& collaboration
    \item \textbf{GitHub Actions:} CI/CD automation
\end{itemize}

\section{Kiến trúc thiết kế}

\subsection{Repository Pattern}
Lớp trung gian giữa Business Logic và Data Access:

\begin{itemize}
    \item Tách biệt database logic từ business logic
    \item Dễ testing (mock repositories)
    \item Dễ thay đổi database implementation
\end{itemize}

\subsection{Service Layer Pattern}
Chứa business logic của ứng dụng:

\begin{itemize}
    \item Validate data
    \item Xử lý business rules
    \item Điều phối repository calls
\end{itemize}

\subsection{DTO (Data Transfer Objects)}
Truyền data giữa API và client:

\begin{itemize}
    \item Tránh truyền entities trực tiếp (security risk)
    \item Kiểm soát fields được expose
    \item Validation dữ liệu input từ client
\end{itemize}

\section{Thống kê mã nguồn}
Dựa trên kho lưu trữ GitHub tại \url{https://github.com/Sunlaii/Si_Sap_InteractHub-main}:

\begin{center}
\begin{tabular}{|l|r|}
\hline
\textbf{Công nghệ} & \textbf{Tỷ lệ} \\
\hline
TypeScript & 68.9\% \\
C\# & 27.4\% \\
TSQL & 2.1\% \\
CSS & 1.4\% \\
Khác & 0.2\% \\
\hline
\end{tabular}
\end{center}

Điều này cho thấy phần frontend (TypeScript) chiếm phần lớn, phản ánh sự phức tạp của giao diện người dùng trong một ứng dụng mạng xã hội.
